\subsection{Monte Carlo Foundations}
% Random sampling, estimators, variance, confidence intervals, law of large numbers in computation.

Monte Carlo methods turn integrals into averages of simulated random variables. This is especially useful when the integral is high-dimensional or when the target density is known only up to a normalizing constant.

\begin{remark} \label{rem:bayesian-integrals-motivate-simulation} \hlt{Bayesian Integrals as a Motivation}\\
For prior density $\pi(\theta)$ and likelihood $\mathcal{L}(\theta)$, write $Z=\int \mathcal{L}(u)\pi(u)\,du$.\\
Then $\pi(\theta\mid X^n)=Z^{-1}\mathcal{L}(\theta)\pi(\theta)$ and $\mathbb{E}[\theta\mid X^n]=Z^{-1}\int \theta\,\mathcal{L}(\theta)\pi(\theta)\,d\theta$.\\
For $\theta=(\theta_1,\ldots,\theta_k)$, marginalization requires integrating over the remaining coordinates, making direct analytic evaluation impractical in high dimension.
\end{remark}

\begin{method} \label{method:basic-monte-carlo-integration} \hlt{Basic Monte Carlo Integration}
\begin{enumerate}[label=(\roman*)]
\setlength{\itemsep}{0pt}
\item Let $I=\int h(x)f(x)\,dx$, where $f$ is a probability density.
\item For IID draws $X_1,\ldots,X_N\sim f$, define $\widehat{I}_N=N^{-1}\sum_{i=1}^{N}h(X_i)$; then $\widehat{I}_N\xrightarrow{p}I$ by the law of large numbers.
\item With $Y_i=h(X_i)$, set $s^2=(N-1)^{-1}\sum_{i=1}^{N}(Y_i-\widehat{I}_N)^2$ and $\widehat{\operatorname{se}}=s/\sqrt{N}$.
\item An approximate $1-\alpha$ Monte Carlo interval is $\widehat{I}_N\pm z_{\alpha/2}\widehat{\operatorname{se}}$.
\end{enumerate}
\end{method}

\begin{example} \label{ex:uniform-monte-carlo-integral} \hlt{Integrating on an Interval}\\
For $X\sim\operatorname{Uniform}(a,b)$, $I=\int_a^b h(x)\,dx=(b-a)\mathbb{E}[h(X)]$.\\
Given IID draws $X_1,\ldots,X_N\sim\operatorname{Uniform}(a,b)$, estimate $I$ by $\widehat{I}_N=(b-a)N^{-1}\sum_{i=1}^{N}h(X_i)$.
\end{example}

\begin{example} \label{ex:monte-carlo-normal-cdf} \hlt{Estimating a Normal CDF}\\
Let $Z\sim N(0,1)$ and $\Phi(x)=\mathbb{P}(Z\leq x)$. For IID draws $Z_1,\ldots,Z_N\sim N(0,1)$, use $h(s)=\mathbf{1}\{s\leq x\}$ and $\widehat{\Phi}(x)=N^{-1}\sum_{i=1}^{N}\mathbf{1}\{Z_i\leq x\}$. Thus $\widehat{\Phi}(x)$ is the empirical fraction of draws not exceeding $x$.
\end{example}

\begin{example} \label{ex:two-binomial-posterior-simulation} \hlt{Two Binomial Proportions}\\
Suppose $X\sim\operatorname{Binomial}(n,p_1)$ and $Y\sim\operatorname{Binomial}(m,p_2)$ are independent, with independent flat priors on $p_1,p_2$.\\
Then $p_1\mid X\sim\operatorname{Beta}(X+1,n-X+1)$ and $p_2\mid Y\sim\operatorname{Beta}(Y+1,m-Y+1)$.\\
For $\delta=p_2-p_1$, draw $P_1^{(i)}\sim\operatorname{Beta}(X+1,n-X+1)$ and $P_2^{(i)}\sim\operatorname{Beta}(Y+1,m-Y+1)$, set $\delta^{(i)}=P_2^{(i)}-P_1^{(i)}$, and summarize the empirical distribution of $\delta^{(i)}$.
\end{example}

\begin{method} \label{method:posterior-function-simulation} \hlt{Simulating a Posterior Function}
\begin{enumerate}[label=(\roman*)]
\setlength{\itemsep}{0pt}
\item Let $\delta=g(p_1,\ldots,p_k)$ be the posterior target.
\item Draw posterior samples $(P_1^{(i)},\ldots,P_k^{(i)})$ for $i=1,\ldots,N$.
\item Set $\delta^{(i)}=g(P_1^{(i)},\ldots,P_k^{(i)})$, use empirical distribution of $\delta^{(i)}$ to estimate posterior distribution of $\delta$.
\end{enumerate}
\end{method}
